\section{System Architecture and Design}

\subsection{Overall System Architecture}
Figure~\ref{fig:arch} shows the block diagram. The system follows a layered client--server
architecture with one separately scaled background worker.

\begin{figure}[H]
\centering
\begin{tikzpicture}[node distance=4mm and 6mm]
% clients
\node[ext,minimum width=3.3cm] (android) {Android app};
\node[ext,minimum width=3.3cm,right=of android] (ios) {iPhone app};
\node[ext,minimum width=3.3cm,right=of ios] (web) {Laptop browser\\(web build)};
\node[draw,dashed,rounded corners,fit=(android)(ios)(web),inner sep=6pt,
      label={[font=\scriptsize\bfseries]above:React Native / Expo client (one TypeScript code base)}] (cl) {};
% api
\node[blk,below=13mm of ios,minimum width=12.6cm,minimum height=9mm] (mw) {Middleware: request id, JSON logs, CORS, error envelope, rate limiting};
\node[blk,below=3mm of mw.south west,anchor=north west,minimum width=2.9cm,minimum height=13mm] (sec) {Security\\JWT $\cdot$ RBAC\\policy};
\node[blk,right=2mm of sec,minimum width=2.9cm,minimum height=13mm] (routers) {REST routers\\auth, me, timetables,\\search, admin};
\node[blk,right=2mm of routers,minimum width=3.2cm,minimum height=13mm] (svc) {Services\\context, uploads,\\conflicts, academics};
\node[blk,right=2mm of svc,minimum width=3cm,minimum height=13mm] (qe) {Query planner\\NLP $\rightarrow$ query plan\\$\rightarrow$ QT01--QT23};
\node[draw,thick,rounded corners,fit=(mw)(sec)(qe),inner sep=5pt,
      label={[font=\scriptsize\bfseries,anchor=south west]north west:FastAPI application (stateless, N processes)}] (api) {};
% bottom
\node[store,below=16mm of routers,minimum width=3.6cm,minimum height=11mm,font=\small] (db) {PostgreSQL 16\\31 tables};
\node[blk,below=16mm of qe,minimum width=4cm,minimum height=11mm] (worker) {Background worker\\Extraction pipeline\\(PyMuPDF, Tesseract, Office)};
\node[store,left=10mm of db,minimum width=3cm,minimum height=11mm,font=\small] (files) {Private file\\storage};
\draw[darr] (cl.south) -- node[lbl,right]{HTTPS, JSON (REST /api/v1)} (mw.north -| cl.south);
\draw[darr] (api.south -| db.north) -- node[lbl,right]{SQL (bound params)} (db.north);
\draw[arr] (worker.west) -- node[lbl,above]{claims jobs} node[lbl,below,font=\tiny]{FOR UPDATE SKIP LOCKED} (db.east);
\draw[arr] (api.south -| files.north) -- node[lbl,left]{originals} (files.north);
\draw[arr] (files.south) to[bend right=25] node[lbl,below]{read file} (worker.south);
\end{tikzpicture}
\caption{Overall system architecture (block diagram)}
\label{fig:arch}
\end{figure}

\paragraph{Responsibilities and interactions.}
The \textbf{client} renders dashboards and forms, caches reads with TanStack Query and keeps the
session in the device keystore. It calls the \textbf{API} over HTTPS with a bearer token. The API
authenticates the user, loads roles from the database, resolves the search context (domain,
timetable, time zone, current time), and runs services and the query engine. Uploads are stored
privately and a \texttt{PROCESSING\_JOB} row is inserted in the same transaction. The \textbf{worker}
claims jobs, extracts and validates the timetable, and writes entries back. \textbf{PostgreSQL} is
the single source of truth, and its row locks are also used for synchronisation.

\subsection{Component and Module Description}
\begingroup\small
\begin{longtable}{@{}>{\raggedright\arraybackslash}p{3.3cm}>{\raggedright\arraybackslash}p{4.4cm}>{\raggedright\arraybackslash}p{3.5cm}>{\raggedright\arraybackslash}p{3.9cm}@{}}
\caption{Modules: purpose, inputs, outputs and dependencies}\label{tab:modules}\\
\toprule
\textbf{Module (path)} & \textbf{Purpose} & \textbf{Input / output} & \textbf{Depends on}\\
\midrule
\endfirsthead
\toprule
\textbf{Module (path)} & \textbf{Purpose} & \textbf{Input / output} & \textbf{Depends on}\\
\midrule
\endhead
Auth \& RBAC \newline{\scriptsize\texttt{app/\allowbreak security}} & Login, tokens, password hashing, role$\rightarrow$scope policy, rate limit & Credentials / tokens; request / principal & \texttt{USER\_ACCOUNT}, \texttt{USER\_ROLE}, Argon2, PyJWT\\
API routers \newline{\scriptsize\texttt{app/\allowbreak api/\allowbreak v1}} & HTTP endpoints, request/response validation & JSON, multipart / JSON envelopes & Services, Pydantic\\
Context resolver \newline{\scriptsize\texttt{services/\allowbreak context.py}} & Chooses domain, timetable, time zone, ``now'' on the server & User + hints / context object & Preferences, settings\\
Upload service \newline{\scriptsize\texttt{services/\allowbreak uploads.py}} & Validates file, stores it, creates archive + job & File / timetable id (\textsc{Queued}) & File types, storage, jobs\\
Extraction \newline{\scriptsize\texttt{app/\allowbreak extraction}} & File $\rightarrow$ grid $\rightarrow$ sections, times, cells $\rightarrow$ validated entries & Bytes / \texttt{ExtractionResult} & PyMuPDF, Tesseract, openpyxl, python-docx\\
Worker \newline{\scriptsize\texttt{app/\allowbreak worker.py}} & Claims and runs jobs, heartbeat, retries, status update & Job rows / entries, status, notification & Extraction, DB\\
Query planner \newline{\scriptsize\texttt{app/\allowbreak query}} & Parse, build a typed query plan, resolve entities, validate, compile into templates (specialised QT01--QT22 or composed QT23) & Text / result envelope & Vocabulary, templates, policy\\
Conflicts \newline{\scriptsize\texttt{services/\allowbreak conflicts.py}} & Detect professor/room/batch overlaps & Timetable id / conflict list & Entries\\
Academics \newline{\scriptsize\texttt{services/\allowbreak academics.py}} & Fresher registration, promotion, transitions & Admin commands / history rows & Student history\\
Admin master data & Departments, years, batches, courses, rooms, aliases, users & Forms / records & DB, audit\\
Notifications \& audit & In-app/push notifications; audit trail & Events / rows & Expo push (optional)\\
Mobile app \newline{\scriptsize\texttt{mobile/\allowbreak src}} & Screens, navigation, state, API client & User actions / rendered UI & REST API\\
\bottomrule
\end{longtable}
\endgroup

\subsection{Architecture Selection and Justification}
We compared three options.
\begin{enumerate}
\item \textbf{Microservices} (separate auth, extraction and search services with a message broker):
      independent scaling, but much more deployment and debugging effort for a single-college workload and
      a team of students, and not possible on a free hosting tier.
\item \textbf{Serverless functions}: no server to manage, but cold starts, execution-time limits
      that do not suit OCR, and vendor lock-in.
\item \textbf{Modular monolith + worker} (chosen): one Python service with clearly separated
      layers (API $\rightarrow$ services $\rightarrow$ query/extraction $\rightarrow$ data). The only
      CPU-heavy part, extraction, runs in a separate worker process so it never blocks API requests.
      A PostgreSQL job table replaces a message broker, so job creation is in the \emph{same
      transaction} as the upload record and there is one less moving part.
\end{enumerate}
The API is \textbf{stateless} (REST~\cite{fielding2000}): any process can serve any request, so the
system scales horizontally by adding processes. A \textbf{layered} design keeps the rules
(security, domain separation) in one place each, which makes them testable.

\subsection{Data Flow Diagrams}
\begin{figure}[H]
\centering
\begin{adjustbox}{max width=\textwidth}
\begin{tikzpicture}
\node[proc,minimum size=34mm,font=\small] (sys) {0.0\\Dynamic\\Timetable\\System};
\node[ext,above left=8mm and 26mm of sys] (stu) {Student};
\node[ext,above right=8mm and 26mm of sys] (prof) {Professor};
\node[ext,below left=8mm and 26mm of sys] (adm) {Admin};
\node[ext,below right=8mm and 26mm of sys] (hod) {HOD /\\Principal};
\draw[arr] (stu.east) to[bend left=10] node[lbl,above,sloped]{questions, personal file} (sys.north west);
\draw[arr] (sys.west) to[bend left=10] node[lbl,below,sloped]{schedule, answers} (stu.south);
\draw[arr] (prof.west) to[bend right=10] node[lbl,above,sloped]{questions} (sys.north east);
\draw[arr] (sys.east) to[bend right=10] node[lbl,below,sloped]{schedule, free rooms} (prof.south);
\draw[arr] (adm.east) to[bend right=10] node[lbl,below,sloped]{timetable file, corrections} (sys.south west);
\draw[arr] (sys.west) to[bend right=30] node[lbl,above,sloped]{validation summary} (adm.north);
\draw[arr] (sys.south east) to[bend left=10] node[lbl,below,sloped]{conflicts, room usage} (hod.west);
\end{tikzpicture}
\end{adjustbox}
\caption{DFD level 0 (context diagram)}
\label{fig:dfd0}
\end{figure}

\begin{figure}[H]
\centering
\begin{adjustbox}{max width=\textwidth}
\begin{tikzpicture}
\node[ext] (adm) at (0,0) {Admin};
\node[proc] (p1) at (3.4,0) {1.0\\Receive \&\\archive};
\node[proc] (p2) at (6.9,0) {2.0\\Extract \&\\validate};
\node[proc] (p3) at (10.6,0) {3.0\\Review \&\\activate};
\node[store] (d1) at (5.15,2.3) {D1 Uploaded files};
\node[store] (d3) at (0,-3.4) {D3 Users, roles,\\preferences};
\node[proc] (p4) at (3.4,-3.4) {4.0\\Resolve\\context};
\node[proc] (p5) at (6.9,-3.4) {5.0\\Answer query\\(query plan\\$\rightarrow$ SQL)};
\node[store] (d2) at (10.6,-3.4) {D2 Timetables\\\& entries};
\node[proc] (p6) at (10.6,-6.6) {6.0\\Detect\\conflicts};
\node[ext] (hod) at (14.4,-6.6) {HOD /\\Principal};
\node[ext] (usr) at (6.9,-6.9) {Student /\\Professor};
\draw[arr] (adm) -- node[lbl,above]{file} (p1);
\draw[arr] (p1) -- node[lbl,above]{job} (p2);
\draw[arr] (p1) |- node[lbl,pos=0.75,above]{original} (d1.west);
\draw[arr] (d1.east) -| (p2);
\draw[arr] (p2) -- node[lbl,above]{entries} (p3);
\draw[arr] (adm.north) -- ++(0,3.4) -| node[lbl,pos=0.25,above]{corrections, activation} (p3.north);
\draw[arr] (p3) -- node[lbl,right]{verified data} (d2);
\draw[arr] (d3) -- (p4);
\draw[arr] (p4) -- node[lbl,above]{context} (p5);
\draw[arr] (d2) -- node[lbl,above]{entries} (p5);
\draw[arr] (d2) -- node[lbl,right]{entries} (p6);
\draw[arr] (p6) -- node[lbl,above]{clashes} (hod);
\draw[arr] ([xshift=-5mm]usr.north) -- node[lbl,left]{question} ([xshift=-5mm]p5.south);
\draw[arr] ([xshift=5mm]p5.south) -- node[lbl,right]{answer} ([xshift=5mm]usr.north);
\end{tikzpicture}
\end{adjustbox}
\caption{DFD level 1}
\label{fig:dfd1}
\end{figure}
The level-0 DFD shows the system with its four external entities. The level-1 DFD splits it into the
upload path (1.0--3.0, used by the Admin) and the query path (4.0--6.0, used by everyone). Both
share the timetable store D2. Personal timetables follow the same path but are written to separate
personal tables (not drawn).

\subsection{UML Diagrams}
\subsubsection{Use case diagram}
\begin{figure}[H]
\centering
\begin{adjustbox}{max width=\textwidth}
\begin{tikzpicture}[uc/.style={draw,ellipse,thick,fill=boxfill,font=\scriptsize,align=center,minimum width=33mm,minimum height=7mm,inner sep=1pt},
  actor/.style={font=\small,align=center}]
\node[draw,thick,rounded corners,minimum width=8.6cm,minimum height=12.3cm,label={[font=\small\bfseries]above:Dynamic Timetable}] (b) at (0,0) {};
\foreach \n/\t/\y in {u1/Sign in/5.4, u2/View dashboard \& week/4.5, u3/Ask question (NL search)/3.6, u4/Find free rooms/2.7,
  u5/Manage personal timetable/1.8, u6/View own teaching \& free time/0.9, u7/Review conflicts/0, u8/View floor activity/-0.9,
  u9/Upload official timetable/-1.8, u10/Review \& correct entries/-2.7, u11/Activate timetable/-3.6, u12/Manage master data \& users/-4.5, u13/Promote students/-5.4}
  { \node[uc] (\n) at (0,\y) {\t}; }
\node[actor] (s) at (-6.2,4) {\stick\\\textbf{Student}};
\node[actor] (p) at (-6.2,0.3) {\stick\\\textbf{Professor}};
\node[actor] (h) at (6.2,0.6) {\stick\\\textbf{HOD /}\\\textbf{Principal}};
\node[actor] (a) at (6.2,-3.6) {\stick\\\textbf{Admin}};
\foreach \u in {u1,u2,u3,u4,u5} \draw (s.east) -- (\u.west);
\foreach \u in {u1,u2,u3,u4,u6} \draw (p.east) -- (\u.west);
\foreach \u in {u1,u2,u7,u8} \draw (h.west) -- (\u.east);
\foreach \u in {u1,u9,u10,u11,u12,u13} \draw (a.west) -- (\u.east);
\end{tikzpicture}
\end{adjustbox}
\caption{Use case diagram}
\label{fig:usecase}
\end{figure}
Professors inherit the student-facing use cases. The HOD and Principal are professors with wider scopes,
which is why the system uses many-to-many \texttt{USER\_ROLE} rather than one role per user.

\subsubsection{Class diagram}
\begin{figure}[H]
\centering
\resizebox{\textwidth}{!}{%
\begin{tikzpicture}[cls/.style={draw,thick,rectangle split,rectangle split parts=3,fill=boxfill,
  font=\scriptsize,align=left,text width=40mm,rectangle split part fill={spitblue!25,boxfill,white}}]
\node[cls] (ua) at (0,0) {\textbf{UserAccount}
\nodepart{second}user\_id: UUID\\college\_email: str\\password\_hash: str\\account\_status\\token\_version: int
\nodepart{third}+roles(): set[Role]};
\node[cls] (st) at (-5,-4.2) {\textbf{Student}
\nodepart{second}uid: str\\student\_status
\nodepart{third}+current\_placement()};
\node[cls] (sf) at (0,-4.2) {\textbf{Staff}
\nodepart{second}short\_code: str\\designation
\nodepart{third}};
\node[cls] (it) at (5.4,0) {\textbf{InstitutionalTimetable}
\nodepart{second}timetable\_id: UUID\\processing\_status\\parser\_version\\effective\_from: date\\data\_version: int
\nodepart{third}+activate()\\+conflicts(): list};
\node[cls] (ie) at (5.4,-4.6) {\textbf{InstitutionalEntry}
\nodepart{second}day\_of\_week, start/end\\course/staff/room/batch\\raw\_extracted\_text\\source\_page\\verification\_status
\nodepart{third}+correct(change)};
\node[cls] (job) at (10.6,0) {\textbf{ProcessingJob}
\nodepart{second}status, attempts\\heartbeat\_at
\nodepart{third}+claim()\\+fail(code)};
\node[cls] (qe) at (-5,0) {\textbf{QueryEngine}
\nodepart{second}registry: Q01..Q30\\templates: QT01..QT22
\nodepart{third}+search(ctx, text)\\-classify(), -resolve()};
\node[cls] (b) at (10.6,-4.6) {\textbf{Batch}
\nodepart{second}code (SE-A1)\\parent\_batch\_id
\nodepart{third}+ancestors()};
\draw[thick] (ua.south) -- node[lbl,pos=.85,right]{0..1} (sf.north);
\draw[thick] (ua.south west) -- node[lbl,pos=.85,left]{0..1} (st.north);
\draw[thick,-{Diamond[open,length=3mm]}] (ie.north) -- node[lbl,pos=0.2,right]{1..*} (it.south);
\draw[thick] (job.west) -- node[lbl,above]{1} (it.east);
\draw[thick,dashed,-{Stealth}] (qe.east) -- node[lbl,above]{reads} (ua.west);
\draw[thick] (ie.east) -- node[lbl,above]{*\quad 0..1} (b.west);
\draw[thick] (sf.east) -- node[lbl,above]{0..1\quad *} (ie.west);
\draw[thick] (b.north) .. controls +(0.8,1) and +(1.6,0.2) .. node[lbl,right]{parent} (b.east);
\end{tikzpicture}}
\caption{Simplified class diagram of the core domain}
\label{fig:class}
\end{figure}

\subsubsection{Sequence diagram: answering a question}
\begin{figure}[H]
\centering
\begin{tikzpicture}[
  ll/.style={draw,thick,fill=boxfill,font=\scriptsize,minimum width=19mm,minimum height=6mm},
  msg/.style={-{Stealth},thick,font=\scriptsize},
  ret/.style={-{Stealth},dashed,font=\scriptsize}]
\foreach \x/\n/\t in {0/c/App, 2.6/a/API router, 5.2/s/Security, 7.8/x/Context, 10.4/q/Query engine, 13.0/d/PostgreSQL}{
  \node[ll] (\n) at (\x,0) {\t};
  \draw[dashed,gray] (\x,-0.3) -- (\x,-9.3);}
\draw[msg] (0,-0.9) -- node[above]{POST /search \{query\}} (2.6,-0.9);
\draw[msg] (2.6,-1.6) -- node[above]{verify JWT} (5.2,-1.6);
\draw[msg] (5.2,-2.2) -- node[above]{load user+roles} (13,-2.2);
\draw[ret] (13,-2.7) -- node[above]{principal} (5.2,-2.7);
\draw[msg] (2.6,-3.4) -- node[above]{resolve(domain, timetable, tz)} (7.8,-3.4);
\draw[msg] (7.8,-4.0) -- node[above]{preferences, settings} (13,-4.0);
\draw[ret] (7.8,-4.6) -- node[above]{context} (2.6,-4.6);
\draw[msg] (2.6,-5.3) -- node[above]{search(ctx, text)} (10.4,-5.3);
\draw[thick] (10.4,-5.6) -- ++(0.6,0) -- node[right,font=\scriptsize,align=left]{parse, classify,\\resolve entities,\\validate params} ++(0,-0.9) -- ++(-0.6,0);
\draw[msg] (10.4,-7.0) -- node[above]{template QTnn (bound params)} (13,-7.0);
\draw[ret] (13,-7.6) -- node[above]{rows} (10.4,-7.6);
\draw[ret] (10.4,-8.2) -- node[above]{envelope: status, results, warnings} (2.6,-8.2);
\draw[ret] (2.6,-8.9) -- node[above]{200 / 422 clarification} (0,-8.9);
\end{tikzpicture}
\caption{Sequence diagram for a natural-language search}
\label{fig:seq}
\end{figure}
Identity, roles and the selected timetable are always resolved on the server, so the client cannot
claim another role or timetable. If a parameter is missing (for example a.m./p.m.), the engine
returns \texttt{CLARIFICATION\_REQUIRED} with choices instead of running a template.

\subsubsection{Activity diagram: publishing a new timetable}
\begin{figure}[H]
\centering
\begin{tikzpicture}[node distance=5mm]
\node[fill=black,circle,minimum size=3mm] (s) {};
\node[blk,below=of s] (a1) {Admin uploads file};
\node[dec,below=of a1] (d1) {type, size,\\signature OK?};
\node[blk,right=14mm of d1,fill=red!10] (rej) {Reject\\(nothing stored)};
\node[blk,below=of d1] (a2) {Store privately, archive row, queue job};
\node[blk,below=of a2] (a3) {Worker: extract, parse, validate};
\node[dec,below=of a3] (d2) {usable?};
\node[blk,right=14mm of d2,fill=red!10] (un) {\textsc{Unusable} /\\\textsc{Failed}};
\node[blk,below=of d2] (a4) {\textsc{Needs Review} or \textsc{Ready}};
\node[blk,below=of a4] (a5) {Admin reviews \& corrects (audited)};
\node[dec,below=of a5] (d3) {activate?};
\node[blk,below=of d3] (a6) {Lock settings, set primary, version+1, notify users};
\node[draw,circle,minimum size=4mm,below=of a6,path picture={\fill (path picture bounding box.center) circle (1.2mm);}] (e) {};
\node[blk,left=14mm of d3] (keep) {Stays in archive};
\draw[arr] (s) -- (a1); \draw[arr] (a1) -- (d1); \draw[arr] (d1) -- node[lbl,above]{no} (rej);
\draw[arr] (d1) -- node[lbl,right]{yes} (a2); \draw[arr] (a2) -- (a3); \draw[arr] (a3) -- (d2);
\draw[arr] (d2) -- node[lbl,above]{no} (un); \draw[arr] (d2) -- node[lbl,right]{yes} (a4);
\draw[arr] (a4) -- (a5); \draw[arr] (a5) -- (d3); \draw[arr] (d3) -- node[lbl,right]{yes} (a6); \draw[arr] (a6) -- (e);
\draw[arr] (d3) -- node[lbl,above]{no} (keep);
\end{tikzpicture}
\caption{Activity diagram: upload, review and activation}
\label{fig:activity}
\end{figure}

\subsubsection{State machine diagram: timetable processing status}
\begin{figure}[H]
\centering
\begin{adjustbox}{max width=\textwidth}
\begin{tikzpicture}[st/.style={draw=spitblue,thick,rounded corners=6pt,fill=boxfill,font=\small,minimum width=26mm,minimum height=8mm}]
\node[fill=black,circle,minimum size=3mm] (i) at (-2.4,0) {};
\node[st] (q) at (0,0) {\textsc{Queued}};
\node[st] (p) at (4.2,0) {\textsc{Processing}};
\node[st] (r) at (11,2.4) {\textsc{Ready}};
\node[st] (n) at (11,0) {\textsc{Needs Review}};
\node[st] (u) at (11,-2.4) {\textsc{Unusable}};
\node[st] (f) at (4.2,-2.8) {\textsc{Failed}};
\draw[arr] (i) -- node[lbl,above]{upload} (q);
\draw[arr] (q) -- node[lbl,above]{worker claims} (p);
\draw[arr] (p.north east) -- node[lbl,above,sloped,pos=0.55]{all entries verified} (r.west);
\draw[arr] (p) -- node[lbl,above,pos=0.55]{warnings / unverified entries} (n);
\draw[arr] (p.south east) -- node[lbl,below,sloped,pos=0.55]{no usable entries / $<$25\% usable} (u.west);
\draw[arr] (p) -- node[lbl,right]{error after 3 attempts} (f);
\draw[arr] (p.north west) .. controls +(-0.3,1.4) and +(0.3,1.4) .. node[lbl,above]{retry (crash / stale heartbeat)} (p.north);
\draw[arr] (n.east) to[bend right=50] node[lbl,right]{review all} (r.east);
\draw[arr] (f.west) -| node[lbl,pos=0.25,below]{reprocess} (q.south);
\end{tikzpicture}
\end{adjustbox}
\caption{State machine of a timetable archive (\textsc{Ready} and \textsc{Needs Review} may be activated)}
\label{fig:state}
\end{figure}

\subsubsection{Deployment diagram}
\begin{figure}[H]
\centering
\begin{adjustbox}{max width=\textwidth}
\begin{tikzpicture}[node/.style={draw,thick,fill=gray!8,align=center,font=\small,minimum height=10mm,inner sep=5pt}]
\node[node] (ph) at (0,0) {\textless\textless device\textgreater\textgreater\\Phone (Android / iOS)\\browser or Expo app};
\node[node] (lp) at (0,-2.6) {\textless\textless device\textgreater\textgreater\\Laptop browser};
\node[node,fill=white,font=\scriptsize] (tls) at (5.3,-1.3) {HTTPS / TLS\\termination};
\node[node,fill=boxfill] (svc) at (9.2,0) {\textless\textless container\textgreater\textgreater\\Docker image: FastAPI\\+ worker thread\\+ static web app};
\node[node,fill=boxfill] (pg) at (9.2,-2.7) {\textless\textless database\textgreater\textgreater\\Managed PostgreSQL 16};
\node[draw,thick,fit=(tls)(svc)(pg),inner sep=8pt,label={[font=\small\bfseries]above:Render cloud (free tier)}] (cloud) {};
\node[node,font=\scriptsize] (gh) at (14.3,-1.3) {GitHub repo\\(push to branch\\$\Rightarrow$ auto deploy)};
\draw[darr] (ph.east) -- node[lbl,above,sloped]{HTTPS} (tls.west);
\draw[darr] (lp.east) -- node[lbl,below,sloped]{HTTPS} (tls.west);
\draw[darr] (tls.east) -- (svc.west);
\draw[darr] (svc) -- node[lbl,right]{private network, TCP 5432} (pg);
\draw[arr] (gh.west) -- (cloud.east);
\end{tikzpicture}
\end{adjustbox}
\caption{Deployment diagram of the free online demo}
\label{fig:deploy}
\end{figure}
For the demo, the worker runs as a thread inside the web container, because the free tier allows one
service. In production the API and the worker run as separate processes or containers
(\texttt{docker-compose.yml}).

\subsection{Database Design}
\subsubsection{ER diagram}
\begin{figure}[H]
\centering
\begin{adjustbox}{max width=\textwidth}
\begin{tikzpicture}[ent/.style={draw=spitblue,thick,fill=boxfill,font=\scriptsize\bfseries,minimum width=27mm,minimum height=7mm,align=center},
  rel/.style={draw,thick,diamond,aspect=2,fill=yellow!15,font=\tiny,inner sep=1pt,align=center}]
\node[ent] (ua) at (0,0) {USER\_ACCOUNT};
\node[ent,minimum width=16mm] (role) at (-6.2,0) {ROLE};
\node[rel] (has) at (-3.4,0) {USER\_ROLE};
\node[ent] (stu) at (-2.5,-2.2) {STUDENT};
\node[ent] (stf) at (2.5,-2.2) {STAFF};
\node[ent] (sah) at (-2.5,-4.4) {STUDENT\_ACADEMIC\\\_HISTORY};
\node[ent] (batch) at (-2.5,-6.8) {BATCH};
\node[ent] (dept) at (-6.6,-5.6) {DEPARTMENT};
\node[ent] (ay) at (-6.6,-3.4) {ACADEMIC\_YEAR};
\node[ent] (itt) at (5.2,-4.4) {INSTITUTIONAL\\\_TIMETABLE};
\node[ent] (ite) at (2.5,-6.8) {INSTITUTIONAL\\\_TIMETABLE\_ENTRY};
\node[ent] (room) at (2.5,-9) {ROOM};
\node[ent] (course) at (-2.5,-9) {COURSE};
\node[ent] (set) at (8.6,-4.4) {SETTINGS\\(single row)};
\node[ent] (ptt) at (5.2,0) {PERSONAL\\\_TIMETABLE};
\node[ent] (pte) at (8.6,0) {PERSONAL\\\_ENTRY};
\node[ent] (job) at (8.6,-6.8) {PROCESSING\\\_JOB};
\node[ent] (cor) at (6,-9) {ENTRY\\\_CORRECTION};
\draw (role) -- node[lbl,above]{M} (has) -- node[lbl,above]{N} (ua);
\draw (ua) -- node[lbl,left]{1 : 0..1} (stu);
\draw (ua) -- node[lbl,right]{1 : 0..1} (stf);
\draw (stu) -- node[lbl,left]{1 : N} (sah);
\draw (sah) -- node[lbl,left]{N : 1} (batch);
\draw (sah) -- node[lbl,above,sloped]{N : 1} (ay);
\draw (batch) -- node[lbl,above,sloped]{N : 1} (dept);
\draw (batch) .. controls +(-1.6,-0.9) and +(-1.6,0.9) .. node[lbl,left]{parent} (batch);
\draw (itt) -- node[lbl,above,sloped]{1 : N} (ite);
\draw (ite) -- node[lbl,above]{N : 1} (batch);
\draw (ite) -- node[lbl,right]{N : 1} (room);
\draw (ite) -- node[lbl,above,sloped]{N : 1} (course);
\draw (ite) -- node[lbl,right]{N : 1} (stf);
\draw (set) -- node[lbl,above]{0..1} (itt);
\draw (ua) -- node[lbl,above]{1 : N owns} (ptt);
\draw (ptt) -- node[lbl,above]{1 : N} (pte);
\draw (job) -- node[lbl,above,sloped]{N : 1} (itt);
\draw (cor) -- node[lbl,above,sloped]{N : 1} (ite);
\end{tikzpicture}
\end{adjustbox}
\caption{ER diagram (main entities; personal sections, notifications, tokens and audit omitted)}
\label{fig:er}
\end{figure}

\subsubsection{Relational schema, keys and constraints}
The database has 31 tables, created by three Alembic migrations. Table~\ref{tab:schema} lists the most
important tables. Primary keys are UUIDs for user-facing objects (non-guessable, which prevents ID
enumeration) and \texttt{BIGINT} identities for master data.

\begingroup\footnotesize
\begin{longtable}{@{}>{\raggedright\arraybackslash}p{4.0cm}>{\raggedright\arraybackslash}p{5.6cm}>{\raggedright\arraybackslash}p{5.6cm}@{}}
\caption{Relational schema of the main tables}\label{tab:schema}\\
\toprule
\textbf{Table (PK)} & \textbf{Main attributes} & \textbf{FKs and constraints}\\
\midrule
\endfirsthead
\toprule
\textbf{Table (PK)} & \textbf{Main attributes} & \textbf{FKs and constraints}\\
\midrule
\endhead
\texttt{user\_\allowbreak{}account}\newline(user\_id) & college\_email, display\_name, password\_hash, account\_status, token\_version & UNIQUE(college\_email)\\
\texttt{user\_\allowbreak{}role}\newline(user\_id, role\_id) & -- & FK user, FK role (M:N bridge)\\
\texttt{student}\newline(student\_id) & uid, student\_status & FK user (UNIQUE), UNIQUE(uid)\\
\texttt{staff}\newline(staff\_id) & short\_code, designation & FK user (UNIQUE), FK department\\
\texttt{student\_\allowbreak{}academic\_\allowbreak{}history}\newline(history\_id) & year\_of\_study, semester, status, effective\_from/to & FK student, academic\_year, department, batch; \textbf{partial UNIQUE(student\_id) WHERE effective\_to IS NULL}\\
\texttt{batch}\newline(batch\_id) & code (SE-A1), cohort, division, program\_year & FK department, self-FK parent\_batch\_id; UNIQUE(dept, code)\\
\texttt{institutional\_\allowbreak{}timetable}\newline(timetable\_id) & title, file\_sha256, file\_format, processing\_status, parser\_version, effective\_from, data\_version & FK academic\_year, department, uploader; partial UNIQUE(uploader, client\_request\_id)\\
\texttt{institutional\_\allowbreak{}timetable\_\allowbreak{}entry}\newline(entry\_id) & entry\_kind, day\_of\_week, class\_date, start\_time, end\_time, raw labels, raw\_extracted\_text, source\_page, confidence\_score, verification\_status, time\_uncertain & FK timetable, section, course, batch, staff, room; CHECK day 1--7; CHECK end$>$start; CHECK status $\in$ \{VERIFIED, UNVERIFIED, INCOMPLETE, REJECTED\}; CHECK 0$\le$confidence$\le$1\\
\texttt{institutional\_\allowbreak{}timetable\_\allowbreak{}settings}\newline(settings\_id) & default\_timetable\_id, version, updated\_by & \textbf{CHECK settings\_id = 1} (single row)\\
\texttt{personal\_\allowbreak{}timetable}\newline(timetable\_id) & owner\_user\_id, status, \ldots & UNIQUE(timetable\_id, owner\_user\_id) for composite FKs\\
\texttt{user\_\allowbreak{}preference}\newline(user\_id) & selection\_mode, personal\_primary\_timetable\_id, last\_active\_domain, last institutional/personal timetable & \textbf{composite FKs (timetable\_id, owner\_user\_id)}: a user can only point to their own personal timetable\\
\texttt{processing\_\allowbreak{}job}\newline(job\_id) & domain, status, attempts, heartbeat\_at & CHECK exactly one target matching the domain\\
\texttt{entry\_\allowbreak{}correction}\newline(correction\_id) & domain, before/after JSON, reason, corrected\_by, corrected\_at & FK entry\\
\texttt{audit\_\allowbreak{}event}\newline(event\_id) & actor, event\_type, target\_type/id, request\_id, occurred\_at & append-only\\
\bottomrule
\end{longtable}
\endgroup

\paragraph{Normalisation.}
The schema is in third normal form: course, staff, room and batch details are stored once and
referenced by FK; many-to-many relations (users--roles) use bridge tables; academic placement is a
separate history table instead of repeating groups in \texttt{student}. One deliberate
\emph{denormalisation} is that each entry also keeps the raw labels exactly as printed (e.g.\
\texttt{KKD}, \texttt{606-4}). This is provenance data rather than derived data: it is needed for
review and stays correct when master data changes.

\paragraph{Indexing strategy.}
Every query template filters by the selected timetable first, so composite B-tree indexes start with
\texttt{timetable\_id}:
\texttt{(timetable\_id, day\_of\_week, start\_time)} for day views, ``now'' and ``next''; \texttt{(timetable\_id,
class\_date, start\_time)} for dated entries; and \texttt{(timetable\_id, staff\_id)}, \texttt{(…, room\_id)},
\texttt{(…, batch\_id)} and \texttt{(…, course\_id)} for professor, room, batch and course searches.
Partial unique indexes enforce business rules (one open placement per student, one active
academic year, idempotent uploads). Status indexes support the worker's job claim.

\paragraph{Sample records.}
Table~\ref{tab:sample} shows real rows extracted from page 1 of the college PDF (SE Division A, Monday).

\begin{table}[H]
\centering
\caption{Sample entries extracted from the real timetable (page 1, Monday)}
\label{tab:sample}
\small
\begin{tabular}{@{}cccllllc@{}}
\toprule
\textbf{Day} & \textbf{Start} & \textbf{End} & \textbf{Course} & \textbf{Staff} & \textbf{Room} & \textbf{Batch} & \textbf{Status}\\
\midrule
1 & 09:00 & 10:00 & LLC & -- & -- & SE-A & VERIFIED\\
1 & 10:00 & 11:00 & DBMS & KKD & 508 & SE-A & VERIFIED\\
1 & 11:15 & 12:15 & DSGT & AAG & 508 & SE-A & VERIFIED\\
1 & 12:15 & 13:15 & DS & AVN & 508 & SE-A & VERIFIED\\
1 & 14:15 & 16:15 & COA Lab & JS & 606-4 & SE-A1 & VERIFIED\\
1 & 14:15 & 16:15 & COA Lab & AVS & 607-B & SE-A2 & VERIFIED\\
\bottomrule
\end{tabular}
\end{table}

\subsection{Security and Error Handling Design}
\begin{table}[H]
\centering
\caption{Threats and controls}
\label{tab:security}
\small
\begin{tabularx}{\textwidth}{@{}p{3.6cm}X@{}}
\toprule
\textbf{Threat} & \textbf{Control in the design}\\
\midrule
Password theft & Argon2id hashing~\cite{biryukov2016}; generic login error; per-IP and per-e-mail rate limit\\
Stolen refresh token & Rotation on every refresh; reuse of an old token revokes the whole token family\\
Viewing policy & The official timetable is public, so every signed-in user may view and ask about all of it; roles only decide what ``my'' means and who may change data (Admin)\\
Role escalation & Roles read from \texttt{USER\_ROLE} on every request; role change bumps \texttt{token\_version}\\
IDOR on personal data & Owner check in the query itself; non-owners get 404 (no existence leak) and an \texttt{ACCESS\_DENIED} audit row\\
SQL / prompt injection & No SQL built from text; query plans use a closed vocabulary of operations and filters; typed, bound parameters; unknown parameters rejected\\
Malicious upload & Extension allow-list \emph{and} magic-byte signature check; size/page limits; private storage; parsing isolated in the worker; LibreOffice with isolated profile and timeout\\
Account takeover via reset & Reset codes never returned in API responses or logs in production (bug P0-1 found and fixed in our audit)\\
Information leakage & Uniform error envelope with request id; no stack traces or SQL in responses; no tokens or free text in logs\\
Secrets in the app & Only the public API URL is in the bundle; tokens in iOS Keychain / Android Keystore (SecureStore)\\
\bottomrule
\end{tabularx}
\end{table}
\paragraph{Error handling.}
Every error returns the same JSON envelope (\texttt{code}, \texttt{message}, \texttt{request\_id}).
Validation errors return 422 with field details. Query results use application statuses
(\texttt{OK}, \texttt{NO\_MATCH}, \texttt{CLARIFICATION\_REQUIRED}, \texttt{UNSUPPORTED\_INTENT},
\texttt{DATA\_UNVERIFIED}), so the app can show a helpful message instead of a generic failure. The
worker separates deterministic failures (unreadable file: no retry) from unexpected ones (retried up
to three times). Structured JSON logs carry the request id for tracing.

\newpage
